% =====================================================================
%  PHAN MO DAU.
%
%  ⚠ DINH VI PHAI CHINH XAC. Fedrecheski et al. CO dem messages va fragments
%  (Bang I cua ho), nen KHONG duoc viet ``chua ai dem''. Cai ho khong dem la so
%  LUOT block-wise, va ly do la o co dien EDHOC chi can 4 manh nen block-wise
%  chua bao gio kich hoat. Che do do cua ho khong cham toi nguong do.
%
%  ⚠ x1.44 la cau hinh DTLS RPK KHONG xac thuc hai chieu. Cau hinh chung thu +
%  hai chieu la x2.79. Trich tran con so ma khong kem cau hinh la sai.
%
%  ⚠ KHONG GO SO. Xem paper/check-no-typed-numbers.py va paper/check-sync.py
% =====================================================================

\section{Introduction}
\label{sec:intro}

EDHOC~\cite{rfc9528} exists because it is smaller than DTLS~1.3~\cite{rfc9147}; it runs over
CoAP~\cite{rfc7252} and establishes an OSCORE~\cite{rfc8613} security context. That is the
argument for deploying it on constrained links, and it is well supported: measuring both
protocols on constrained hardware over an IEEE~802.15.4 radio, Fedrecheski et
al.\ \cite{fedrecheski2024} report a $\times\citedFedrPktLo$ to $\times\citedFedrPktHi$ reduction in packet sizes.

The same measurement, however, records something the size argument does not predict. Against
DTLS with raw public keys and no mutual authentication, that reduction buys a $\times\citedFedrDur$
improvement in handshake duration; the authors state the point themselves, that ``the
improvement in packet sizes in EDHOC actually translates to a smaller improvement in duration
and energy''. Bytes on the wire and cost on the link are already only loosely coupled, with
classical credentials.

Post-quantum credentials change the size regime rather than merely the sizes. Their Table~I
records \citedFedrEdhocFrags{} fragments for a mutually authenticated EDHOC handshake: each message fits inside
one IEEE~802.15.4 frame, so CoAP block-wise transfer~\cite{rfc7959} never engages and the
handshake completes in the two round trips the EDHOC-over-CoAP profile advertises. The profile
itself makes that advertisement conditional, noting that the advantage ``can be lost'' in
combination with block-wise transfers, and that the two-round-trip minimum depends on
message\_3 remaining ``relatively small \ldots within target MTU sizes'' \cite{rfc9668}.

Independent analyses of EDHOC's security~\cite{perez2024} and later extensions such as
efficient rekeying~\cite{lopezgomez2026} take that size advantage as given.

A post-quantum key encapsulation key does not remain within those sizes. The smallest
standardised choice, ML-KEM-512, has an encapsulation key of \numKemFiveOneTwo~B against a
frame payload of \numFramePayload~B. Post-quantum EDHOC is by now implemented and measured on
real hardware \cite{fraile2025}, over 6LoWPAN on BLE, establishing that it runs. Whether it is
still preferable to the protocol it was designed to replace is a separate question, and it is
the one asked here.

This letter reports two results. First, the condition the profile names is not merely
approached but crossed: block-wise engages in message\_1 alone, at every NIST parameter set
and whatever authentication is chosen, so the exchange count grows with size while DTLS, which
fragments at its own handshake layer and sends each flight back-to-back, does not. Second, and
independently of any post-quantum consideration, the DTLS implementations one would compare
against do not all run on such a link at all: of three widely used libraries, one cannot
complete a handshake at an IEEE~802.15.4 frame payload with any credential size, classical
ones included, and the three are bounded by different quantities. DTLS implementations have
been compared before, notably by state fuzzing across thirteen of
them~\cite{fiterau2020}, but for state-machine conformance rather than behaviour under a
constrained MTU; that work states explicitly that its harness ``does not employ reordering and
fragmentation''.

Neither result concerns cryptographic strength, and neither requires a post-quantum library to
reproduce: the argument is about message size, so credentials are padded to post-quantum
dimensions rather than computed. All code and measurements are released.
